% Bursting and budding: where the advantage lives --- plan \S3.3, "\S5".
% Largely new.  This section discharges O1-O4 (plan \S1.3): the evolutionary
% framing, the selection argument over the phenotypes of \cref{p:sec:reach},
% the placement of the flooding result relative to the spectrum, and the
% "which pathogens sit where" scope \ChDist{} asked for.
% Source material: 03_fixed_removal.tex:88 (the batch-count argument),
% :275 (the \yp{} argument).
% NOTHING from \cref{p:sec:reach} is rebuilt here; every mechanism is cited.

\section{Bursting and budding: where the advantage lives}
\label{path:sec:selection}

\subsection{What is settled, and what was left open}
\label{path:sec:settled}

\ChPopcap{} settles the comparison under linear clearance, and it does so
twice. First, at fixed total per-cell output $V_\infty$, replacing constant
budding by bursting leaves the basic reproduction number exactly unchanged
(\cref{p:thm:R0inv}): under linear clearance the release schedule is
invisible to $R_0$. Second, the schedule is nevertheless not invisible to
establishment. \Cref{p:thm:flood} shows that bursting strictly lowers the
extinction probability of a small inoculum, relative to matched-mean
budding, if and only if the flooding parameter
\begin{equation}
L:=a(1-b)=a-\frac{\mu}{\lambda}
\label{path:eq:L}
\end{equation}
exceeds one, equivalently $1/\delta<1/\lambda+1/\mu$
(\cref{p:eq:floodcrit}); at $L=1$ the two offspring laws coincide exactly
and every functional of them ties (\cref{p:prop:L1}).

The letter $R_0$ in \ChPop{} denotes the expected number of secondary
infections produced by one infected cell in a wholly susceptible population
under linear clearance, in the sense of \cite{diekmann1990definition}; a
theorem there proves that quantity invariant. The quantity of
\cref{path:eq:R0} is not that: it is the expected number of \emph{surviving}
released agents per infected cell when the clearance is a threshold rather
than a rate. Under linear clearance at infection probability $q$ the two
agree up to the thinning factor, since $\mathcal{R}(0)=V_\infty$ and the
mean offspring number in \ChPop{} is $qV_\infty$. Under threshold removal
they part company, and hence the different letter. Reading
$\mathcal{R}(\vartheta)$ as $R_0$ would put this chapter in contradiction
with \cref{p:thm:R0inv}, and there is none.

\Cref{p:rem:flood-scope} records that both results in \ChPop{} compare
branching structures at fixed offspring mean, and therefore say nothing
about nonlinear effects. Fixed removal is exactly such an effect --- a
threshold is not a rate --- so the comparison can be run again through it.
It does not come out the same way.

\subsection{The reversal}
\label{path:sec:reversal}

Two arguments will appear, and they are not the same argument. The
distributional one holds the number of deliveries fixed at one and compares
two laws for that delivery, matched at mean $V_\infty$; its verdict depends
on $L$. The batch-count one holds the law fixed and asks how the same total
fares when split across more deliveries, each paying a toll $\vartheta$;
its verdict is unconditional concentration. The reversal below is the
first. The spectrum figure is the second.

\begin{proposition}[The flooding criterion reverses under threshold removal]
\label{path:prop:reversal}
Let the budding comparator of \cref{p:sec:flooding} release a total that is
geometric on $\{0,1,2,\dots\}$ with the same mean $V_\infty$, and write
$\sigma=V_\infty/(1+V_\infty)$. Then under an integer removal budget
$\vartheta$,
\begin{equation}
\mathcal{R}_{\mathrm{burst}}(\vartheta)=V_\infty\,a^{-\vartheta},
\qquad
\mathcal{R}_{\mathrm{bud}}(\vartheta)=V_\infty\,\sigma^{\vartheta},
\label{path:eq:reversal}
\end{equation}
so bursting exceeds matched-mean budding if and only if $1/a>\sigma$, that
is if and only if $V_\infty(a-1)<1$, that is if and only if $L<1$ --- the
exact complement of \cref{p:thm:flood}. At $L=1$ the two laws coincide by
\cref{p:prop:L1} and every functional ties.
\end{proposition}

\begin{proof}
Both expressions in \cref{path:eq:reversal} follow from the memorylessness
argument of \cref{path:prop:ratio}, applied to a geometric law on
$\{1,2,\dots\}$ with ratio $1/a$ in the first case and to one on
$\{0,1,2,\dots\}$ with ratio $\sigma$ in the second. For the criterion,
take \cref{path:prop:ratio} at $\vartheta=1$ against the tail
\cref{path:eq:tail}: $V_\infty/a=\ex{r}-\pr{r\ge1}=V_\infty-(1-b)$, whence
$V_\infty=a(1-b)/(a-1)=L/(a-1)$. Therefore
\[
\frac1a>\frac{V_\infty}{1+V_\infty}
\iff 1+V_\infty>aV_\infty
\iff V_\infty(a-1)<1
\iff L<1. \qedhere
\]
\end{proof}

The result is a reversal rather than a contradiction, and the reason is a
single fact about variance. At matched mean, extinction probability rewards
\emph{low} offspring variance, whereas $\ex{\max(0,r-\vartheta)}$ is a convex
functional of the offspring count and so rewards \emph{high} variance.
\Cref{p:eq:varorder} gives the variance ordering in terms of the same
parameter,
$\var{\,\cdot\,}_{\mathrm{bud}}-\var{\,\cdot\,}_{\mathrm{burst}}
=2q^{2}V_\infty(L-1)/(a-1)$. The sign of $L-1$ therefore drives both
criteria, and drives them opposite ways: where bursting is the safer bet for
establishment it is the worse bet against a threshold, and conversely. Both
criteria turn on the offspring count being integer-valued. Only an integer
count lets a threshold bite, and the same feature governs establishment
\cite{williams2024reproduction}.

\begin{table}[tbp]
\centering
\small
\renewcommand{\arraystretch}{1.1}
\begin{tabular}{@{}lccccll@{}}
\toprule
& & & & & \multicolumn{2}{c@{}}{winner under} \\
\cmidrule(l){6-7}
$(\lambda,\mu,\delta)$ & $L$ & $\vartheta$
& $\mathcal{R}_{\mathrm{burst}}$ & $\mathcal{R}_{\mathrm{bud}}$
& linear clearance & threshold removal \\
\midrule
$(1,0,0.1)$    & $1.100$ & $5$  & $6.830$ & $7.120$ & bursting & budding \\
$(1,0.5,1/3)$  & $1.000$ & $5$  & $0.263$ & $0.263$ & tie      & tie \\
$(1,0.9,0.1)$  & $0.416$ & $5$  & $0.333$ & $0.078$ & budding  & bursting \\
$(1,0.2,0.05)$ & $0.862$ & $10$ & $7.692$ & $7.011$ & budding  & bursting \\
$(1,0.9,0.02)$ & $0.207$ & $5$  & $1.164$ & $0.241$ & budding  & bursting \\
$(1,0.95,0.5)$ & $1.017$ & $3$  & $0.138$ & $0.142$ & bursting & budding \\
\bottomrule
\end{tabular}
\caption{The reversal in six regimes. The last two columns name which
strategy wins under linear clearance and which under threshold removal, and
they disagree in every row except the tie. The middle row is the boundary
$L=1$, where the bursting offspring law is exactly geometric
(\cref{p:prop:L1}) and the two ratios agree to the digits shown. The
fourth row is strongly supercritical on both sides, so the reversal is not
an artefact of one of them being near extinction.}
\label{path:tab:reversal}
\end{table}

\begin{figure}[htbp]
\centering
\includegraphics[width=0.74\textwidth]{flooding_reversal.pdf}
\caption{The sign change, drawn. The gap
$\mathcal{R}_{\mathrm{burst}}-\mathcal{R}_{\mathrm{bud}}$ against $\delta$
at $\lambda=1$, $\mu=0.5$, for three removal budgets. All three curves cross
zero at the same place, the flooding boundary
$\delta_{*}=\lambda\mu/(\lambda+\mu)=1/3$ of \cref{p:eq:floodcrit}, and they
cross it in the direction opposite to \cref{p:thm:flood}: bursting wins to
the \emph{left}, where $L<1$. Raising the budget deepens the advantage on
the left without moving the boundary, because the boundary is a property of
the two laws and not of the toll.}
\label{path:fig:reversal}
\end{figure}

\Cref{path:tab:reversal,path:fig:reversal} exhibit the reversal. The
boundary is unmoved by $\vartheta$. Carry that forward: the removal budget
sets how much the winner wins by, and the rate constants alone set who
wins.

If the toll is paid \emph{per delivery}, a cell that delivers a total yield
$V_\infty$ in $n$ instalments retains $\max(0,V_\infty-n\vartheta)$, against
a bursting cell's $\max(0,V_\infty-\vartheta)$. Since $n\ge1$, concentration
wins for every $n$, every $\vartheta$ and every $L$: there is no threshold
and no reversal in this version. That is the batch-count argument of
\cref{path:sec:playingfield}, not \cref{path:prop:reversal}.
\label{path:sec:twoarguments}

\begin{remark}[Scope of the comparison]
\label{path:rem:scope}
The two arguments answer different questions and neither subsumes the other.
The distributional one holds the number of deliveries fixed at one and asks
which \emph{law} for that delivery does better, and its answer depends on
$L$. The batch-count one holds the law fixed and asks how the same total
fares when split across more deliveries, and its answer does not depend on
$L$ at all. Conflating them produces the impression that flooding is
unconditionally advantageous under depletion --- true of the second
argument, false of the first.
\end{remark}

\subsection{Placement on the spectrum}
\label{path:sec:spectrum}

\Cref{p:sec:reach} lays the release mechanisms out as a spectrum
(\cref{p:fig:spectrum}). The catalogue is built there and is not rebuilt
here. What a per-delivery toll does at each station follows from the batch
count without a new derivation: the surviving output is set by how many
times $\vartheta$ is paid.

\begin{figure}[!h]
\centering
\begin{tikzpicture}[>=Stealth, font=\footnotesize]
\draw[->, line width=0.8pt, color=inksoft] (0,0) -- (12.4,0);
\node[below, font=\scriptsize\itshape, text=inksoft] at (0.15,-0.15) {budding};
\node[below, font=\scriptsize\itshape, text=inksoft] at (12.15,-0.15) {bursting};
\foreach \x/\lab/\toll in {%
  1.3/{constant $p$}/{pays $n\vartheta$, $n$ large},%
  3.7/{stock-and-export}/{pays per export},%
  6.2/{reset / partial $\varphi$}/{pays $\sim\vartheta/\varphi$},%
  8.8/{full clear}/{pays once per dump},%
  11.2/{absorbing BDC}/{pays once}}
{\draw[na1] (\x,-0.08) -- (\x,0.08);
 \node[above=0.22cm, align=center, font=\scriptsize] at (\x,0) {\lab};
 \node[below=0.38cm, align=center, font=\scriptsize, text=inksoft] at (\x,0) {\toll};}
\end{tikzpicture}
\caption{Per-delivery toll along the budding--bursting spectrum of
\cref{p:fig:spectrum}. Moving right, fewer deliveries are made and a
smaller total is paid. The ordering is the batch count, not the
distributional reversal.}
\label{path:fig:tollspectrum}
\end{figure}

Under a per-delivery toll, selection on release phenotypes is therefore
monotone toward fewer, larger deliveries
(\cref{path:fig:tollspectrum,p:fig:spectrum}). That is a different
statement from \cref{path:prop:reversal}, which compares two laws at one
delivery each and can go either way. It is also what makes depletion a
candidate selective pressure at all, in the sense in which nested
within-host models pose the question \cite{gilchrist2006evolution}: it
supplies a gradient rather than merely a difference.

That gradient does not run unopposed. It runs against the one \ChPop{}
identifies, since concentrating release lengthens the generation time and so
lowers the deterministic growth rate \cite{wallinga2007generation}, and the
phage lysis-timing literature poses the same trade-off in its own terms
\cite{wang2006lysis}; how a lytic strategy can be evolutionarily competitive
at all under that tension is asked directly \cite{komarova2007viral}. The
contribution here is to name the third quantity in the trade-off and to give
it a closed form.

Self-excitation (\cref{p:sec:selfexcite}) is the one station where the toll
itself moves. \Cref{path:rem:reset} assumes the killer field has recovered
between deliveries; a self-exciting mechanism makes deliveries cluster in
time, so successive dumps meet a field that has not recovered and pay less
than $\vartheta$ each. Whether that is enough to overcome the batch-count
penalty of making several deliveries instead of one is not settled here ---
it depends on the recovery timescale, which the model does not carry --- and
it is recorded as an open problem in \cref{path:sec:open}. It is worth
naming because it is the only mechanism in the catalogue for which the two
arguments of \cref{path:rem:scope} pull in opposite directions.

This is a statement about which release phenotypes a depletable defence
favours, not about which pathogens possess which phenotypes, and it holds
at the level of one cell's output. It inherits the restriction that
\cref{p:thm:flood} and \cite{pearson2011stochastic} carry: what is being
described is the early phase \cite{hataye2019principles}.

\subsection{\texorpdfstring{\yp{}}{Y. pestis}}
\label{path:sec:ypestis}

One reading takes the model of \cref{path:sec:removal} in the context of
\yp{} and neutrophils: the target cells are macrophages, the killer cells
are neutrophils. \yp{} does replicate within macrophages
\cite{PujolBliska2003,KeEtAl2013}, subverting the phagosome rather than
being destroyed by it \cite{ConnorEtAl2018,PujolBliska2005}, and it
encounters those cells early in the pulmonary route
\cite{BosioEtAl2005}. The injected effectors of
\cref{path:sec:chainedruptures} act on the same cells, driving or
withholding apoptosis
\cite{MonackEtAl1997,ZaubermanEtAl2006,PetersEtAl2013}. Both actors of the
model are present. The formulae themselves do not depend on that reading.

\begin{table}[!htbp]
\centering
\small
\begin{tabular}{@{}ll@{}}
\toprule
Model & Reading \\
\midrule
replicators $X_t$ & intracellular pathogen in one macrophage \\
granules $Q_t$ & remaining oxidative store \\
$\varsigma$ & suppression per replicator--granule pair \\
burst size $r$ & load at rupture \\
budget $\vartheta$ & killer cells removed per delivery \\
$\mathcal{R}(\vartheta)$ & expected surviving release \\
\bottomrule
\end{tabular}
\caption{A reading of already-named objects under the \yp{} mapping.
The formulae themselves do not depend on it.}
\label{path:tab:ypestis}
\end{table}

Read that way, the argument is two-phase. While killer-cell density is
roughly constant, the bursting advantage of \cref{path:sec:spectrum}
applies: a macrophage releasing its whole load at once pays one toll where
a budding cell would pay many. If a later rise in killer cells pushes the
toll past $\vartheta^{\ast}=\log V_\infty/\log a$, no intracellular strategy
remains supercritical (\cref{path:cor:threshold}), and the advantage lies
with becoming extracellular. Why not replicate extracellularly from the
beginning is then a separate question: being within macrophages may buy
transport to lymph nodes, and perhaps food itself.

% NEEDS-BIOLOGY: that \yp{} sits at the bursting end of the release
% spectrum, and what intracellular constraint puts it there (bacterial size
% relative to the membrane machinery is the candidate reason given in
% \cref{path:sec:advantages}).
% Candidate sources already copied into references.bib: PujolBliska2003,
% KeEtAl2013, ConnorEtAl2018.
%
% NEEDS-BIOLOGY: that plague infection proceeds in two phases, an early
% intracellular (macrophage) phase followed by a predominantly
% extracellular one, and what the histology shows about the transition.
% Candidate sources already copied into references.bib: BosioEtAl2005,
% PetersEtAl2013, MonackEtAl1997.
%
% NEEDS-BIOLOGY: that neutrophil density at the site of infection rises
% sharply between those two phases, with an order of magnitude if one is
% available; this is the quantity that plays the role of \vartheta and the
% whole two-phase argument turns on it.
% Candidate sources already copied into references.bib: PetersEtAl2013,
% BosioEtAl2005.  No entry in the donor pool states a density.
%
% NEEDS-BIOLOGY: that infected macrophages transport \yp{} to draining
% lymph nodes, and that this is an advantage of the intracellular phase
% sufficient to explain why it comes first.
% Candidate sources already copied into references.bib: BosioEtAl2005,
% PujolBliska2005.
%
% NEEDS-BIOLOGY: whether HIV's observed release clustering is detection
% avoidance, which is the third advantage of \cref{path:sec:advantages} and
% is named there and not pursued.
% Candidate sources already copied into references.bib: conway2018modeling,
% hataye2019principles.

What the mathematics supplies is a single number. At
$(\lambda,\mu,\delta)=(1,0.5,1/3)$ one has $\vartheta^{\ast}=1.71$, so a
killer field able to remove two agents per delivery already forces the
switch; at $(1,0.2,0.05)$ it is $44.13$. If effector density at the site
crosses $\vartheta^{\ast}$ during the infection, the intracellular phase
should end at a level set by the intracellular rate constants and by
nothing else. An analogous reading for \ft{}
\cite{oyston2004tularaemia,carruthers2020stochastic} would use the same
threshold with different constants.

\FloatBarrier
